%
\PassOptionsToPackage{unicode=true}{hyperref} % options for packages loaded elsewhere
\PassOptionsToPackage{hyphens}{url}
%
\documentclass[
]{article}
\usepackage{lmodern}
\usepackage{amssymb,amsmath}
\usepackage{ifxetex,ifluatex}
\ifnum 0\ifxetex 1\fi\ifluatex 1\fi=0 % if pdftex
  \usepackage[T1]{fontenc}
  \usepackage[utf8]{inputenc}
  \usepackage{textcomp} % provides euro and other symbols
\else % if luatex or xelatex
  \usepackage{unicode-math}
  \defaultfontfeatures{Scale=MatchLowercase}
  \defaultfontfeatures[\rmfamily]{Ligatures=TeX,Scale=1}
\fi
% use upquote if available, for straight quotes in verbatim environments
\IfFileExists{upquote.sty}{\usepackage{upquote}}{}
\IfFileExists{microtype.sty}{% use microtype if available
  \usepackage[]{microtype}
  \UseMicrotypeSet[protrusion]{basicmath} % disable protrusion for tt fonts
}{}
\makeatletter
\@ifundefined{KOMAClassName}{% if non-KOMA class
  \IfFileExists{parskip.sty}{%
    \usepackage{parskip}
  }{% else
    \setlength{\parindent}{0pt}
    \setlength{\parskip}{6pt plus 2pt minus 1pt}}
}{% if KOMA class
  \KOMAoptions{parskip=half}}
\makeatother
\usepackage{xcolor}
\IfFileExists{xurl.sty}{\usepackage{xurl}}{} % add URL line breaks if available
\IfFileExists{bookmark.sty}{\usepackage{bookmark}}{\usepackage{hyperref}}
\hypersetup{
  pdfborder={0 0 0},
  breaklinks=true}
\urlstyle{same}  % don't use monospace font for urls
\usepackage{longtable,booktabs}
% Allow footnotes in longtable head/foot
\IfFileExists{footnotehyper.sty}{\usepackage{footnotehyper}}{\usepackage{footnote}}
\makesavenoteenv{longtable}
\usepackage{graphicx,grffile}
\makeatletter
\def\maxwidth{\ifdim\Gin@nat@width>\linewidth\linewidth\else\Gin@nat@width\fi}
\def\maxheight{\ifdim\Gin@nat@height>\textheight\textheight\else\Gin@nat@height\fi}
\makeatother
% Scale images if necessary, so that they will not overflow the page
% margins by default, and it is still possible to overwrite the defaults
% using explicit options in \includegraphics[width, height, ...]{}
\setkeys{Gin}{width=\maxwidth,height=\maxheight,keepaspectratio}
\setlength{\emergencystretch}{3em}  % prevent overfull lines
\providecommand{\tightlist}{%
  \setlength{\itemsep}{0pt}\setlength{\parskip}{0pt}}
\setcounter{secnumdepth}{-2}
% Redefines (sub)paragraphs to behave more like sections
\ifx\paragraph\undefined\else
  \let\oldparagraph\paragraph
  \renewcommand{\paragraph}[1]{\oldparagraph{#1}\mbox{}}
\fi
\ifx\subparagraph\undefined\else
  \let\oldsubparagraph\subparagraph
  \renewcommand{\subparagraph}[1]{\oldsubparagraph{#1}\mbox{}}
\fi

% set default figure placement to htbp
\makeatletter
\def\fps@figure{htbp}
\makeatother


\date{}

\begin{document}

\textbf{Social Finance}

January 2017

\emph{PRELIMINARY DRAFT}\footnote{This is a very preliminary draft. We
  would like to thank Leonardo Becchetti, Robert Lensink, and Ingmar
  Nyman for useful comments and discussions as well as to seminar
  participants at , , Università di Roma Tor Vergata, and . We
  gratefully acknowledge research funding from the Ford Foundation in
  support of this project.}

\textbf{Jonathan Conning\\
}Hunter College and The Graduate Center\\
City University of New York

\textbf{Jonathan M}o\textbf{rduch\\
}New York University

\textbf{Abstract:} We propose a framework for understanding how social
investors' who seek to maximize a combination of private and social
returns from investments in a portfolio of new and established
microfinance institutions. The model takes into account the endogeneity
of loan contract terms as well as the capital structure of the financial
institutions that may emerge to serve target groups of borrowers
differentiated primarily by their levels of initial average net worth,
and how social investments might transform those patterns. We build upon
one of the workhorse models of modern corporate finance (Tirole, 2007)
which features limited liability, multiple layers of moral hazard and
costly monitoring to explain patterns of financial intermediation, and
in our framework, the role and modes of social investment. We pinpoint
the role of the subsidies and guarantees implicit in social investors'
equity and quasi-equity investments and the role they play in attracting
private capital investors and sustaining productivity-enhancing
financial intermediation that might otherwise not have taken place.

\hypertarget{introduction}{%
\section{Introduction}\label{introduction}}

The working modes of philanthropy and business are blurring in two ways.
First, philanthropists are increasingly using methods borrowed from
for-profit investors. The strategies of venture capitalists are being
adapted by ``venture philanthropists'' seeking to fund socially-minded
institutions with promising ideas at early stages of development. They
hope that the ideas will form the basis of sustainable institutions
capable of serving populations at wide scale. As a recent survey of
philanthropy put it:

\begin{quote}
The modern philanthropist dislikes the term charity, preferring to speak
about his social investments. His language is entrepreneurial, sprinkled
with references to metrics, scalability, leverage, and venture
philanthropy (Weinsberg, 2006).
\end{quote}

The changes are expanding to foundations with the increasing favor of
program-related investments over grants (PRIs include loans, equity
shares, recoverable grants, linked deposits, and guarantees); parallel
changes are afoot within government agencies (Marble, 1989).

Oikocredit, for example, a Dutch social investment organization, has
since 1978 been making long-term loans to ``enterprises and institutions
for the empowerment of poor people in developing countries .'' Their
first ``social investment'' in microfinance, for example, was a
\$150,000 loan for 5 years at 9\% interest to Alternativa, a
non-governmental organization providing microcredit to street vendors
and small-scale entrepreneurs in Lima, Peru. Between 1994 and the end of
2004, Oikocredit made a cumulative total of 285 loans or investment
agreements to 161 institutions in 41 countries, totaling \$88 million.
Their net profit, returned to shareholders as a dividend, was 2\%. The
dividend allowed investors a modest return on capital, but they earned
substantially less than alternative investments. Their relative loss is
their implicit charitable contribution (Heinen, 2005).

The second turn is that many socially-minded institutions receiving
funding also increasingly look like businesses. Museums have raised
entry prices and run ambitious shops with mail-order catalogues.
Non-profit schools are hiking tuitions. Environmental groups buy and
sell land. Community development programs work with banks to facilitate
mortgages in low-income neighborhoods. Indeed, some of these
institutions are forming commercially viable entities and competing in
the marketplace on the strength of their balance sheets alone.
Opportunity International, a US-based network of microfinance
institutions, owns a for-profit insurance company. BRAC, the largest NGO
in Bangladesh, owns a for-profit bank. Chicago's Shorebank, the largest
community development bank in the United States, is a holding company
composed of non-profit and for-profit entities. BASIX has adapted the
Shorebank model to South India.

The microfinance sector is home to some of the best-known examples of
the new hybrids, and is in many ways the leading edge. Muhammad Yunus,
the 2006 Nobel Peace Prize winner and founder of the Grameen Bank, a
financial institution for the poor of Bangladesh, created a ``social
business'' funded by social investors.

We tackle the logic and tensions inherent in social finance---the
support, with philanthropic objectives, of nonprofits, social businesses
like the Grameen Bank, and profit-maximizing businesses serving the
poor. The principles behind the new world of philanthropy and social
action have not been well-explored by economists. We develop a theory of
``social finance'' to parallel the modern theory of corporate finance.

Most generally, we assess the justification for the subsidy inherent in
social investment. We go further to describe a framework that describes
how subsidy can be deployed for optimal impact. A focus is on using
funds to build institutions that can achieve scale and attract
additional revenue sources and investments. Leverage is critical, and we
describe two important kinds of leverage. One, having to do with
mobilizing locally-held capital, is often overlooked and is highlighted
below.

Our consideration of effective social investment begins with two ideas.
The first is that effective giving can lead to multiplier effects. It
can ``crowd in'' other resources by catalyzing the investments of
others, increasing the recipient's ability to ``leverage'' its
resources. Leverage (i.e., to be able to obtain funds to use in addition
to an institution's own equity) is critical as it determines an
institution's ability to reach wide scale. Effective social investment
can be used to expand leverage and hence market coverage, holding
constant other social concerns.

The second consideration is that philanthropic giving can change the
incentives for recipients to work hard (as can commercial lending).

In this framework, profitability is neither necessary nor sufficient for
attracting commercial capital. Contrary to some popular claims, earning
profits -- even high profits --will frequently not expand the leverage
of microfinance to take those institutions to wide scale. The presence
of a social investor can however tip the scales for private resources to
pour in. Total output may be enhanced by supporting institutions that
would not exist if profitability was the sole concern. Of course it is
no great surprise that a subsidized investment can expand an activity or
a sector, the claim however is that social investors and philanthropic
capital can expand output and welfare by repairing and expanding
markets.

The framework has a specific implication for the most effective use of
social investment. We describe arguments that push social investment
either ``high'' or ``low'' when directed to poor communities.
Specifically, we show that smart subsidy involves allocations that
either go to institutions serving relatively well-off (but still poor)
customers, where the potential for leverage is maximized, or to
institutions serving the poorest customers. There is less ``bang for the
buck'' in supporting a middle range of institutions that are
self-sufficient with their own resources but for which leverage (and
thus potential scale) is limited.

\hypertarget{two-layers-of-moral-hazard}{%
\subsection{Two layers of moral
hazard}\label{two-layers-of-moral-hazard}}

The analysis focuses on social investment in microfinance institutions.
The particular gap filled by social investors is created by the
unwillingness of most commercial banks to serve the poor. One well-known
difficulty for banks stems from moral hazard on the part of borrowers.
In practice, though, moral hazard is mitigated by microfinance
institutions through the development of technologies, contracts, and
management practices that limit risk. All of this takes effort, but the
high repayment rates achieved by leading practitioners are testimonies
to impressive successes. The essential insight in mitigating moral
hazard is that incentives are improved by exposing people to the
consequences of their actions. In the microfinance context, the denial
of new loans in the case of default (via ``progressive lending'') and
the peer pressure associated with group lending can reduce moral hazard
by adding serious repercussions for bad outcomes (Armendariz de Aghion
and Morduch, 2005).

Given the success in controlling moral hazard, investors ought to be
interested in funding microfinance institutions and expanding their
scale. Investments should be pouring in, and they are starting, but most
investments are made by social investors with explicit philanthropic
objectives.

While social investors are willing to forego the kind of financial
returns that alternative investments could bring, they nonetheless work
hard to preserve their capital. The result is that most money channeled
through microfinance investment funds and via International Financial
Institutions supports ``top-tier,'' mature, financially self-sufficient
institutions (Abrams and von Stauffenberg, 2007). The growing,
not-yet-profitable institutions that are needed to expand the sector are
targeted by less than one fifth of the funds in the MicroBanking
Bulletin's survey of funds (Barres, 2005:53).

One of the chief difficulties for investors---whether driven by
commercial or social ends--is a second layer of moral hazard. The basic
moral hazard problem described above involves the microfinance
institution and the customer. This second problem involves the investor
and the microfinance institution. Once it is the investor's resources
that are ultimately at risk (rather than just the microfinance
institution's resources or the customers'), the incentives are reduced
for the institution to be diligent in managing finances, monitoring
loans, and innovating. The incentives to properly screen potential
customers weaken, as do the incentives to chase down delinquent
borrowers.

The risks are not trivial. Oikocredit, for example, as of late 2004
faced a portfolio at risk (90-day) of 7 percent on its \$67 million
social investment portfolio (Heinen 2005: 33). In a review of 104 loan
contract disbursed before the end of 2002, the average repayment rate
was about 91 percent. Seventy four percent of contracts were found to be
repaid ``almost always on time,'' but ``15 percent require continuous
monitoring to follow overdue payments closely, and another 11 percent
can be classified as bad-debt cases that are persistent in their
default. The latter cases come with high monitoring and legal costs and
are likely to lead to (partial) loan loss in the end.'' The challenges
among the African institutions are greater: ``Only 57 percent of MFIs in
the African sample manage to pay ``on time,'' while more than one out of
four are structurally delayed in payment of invoiced amounts (Heinen
2005: 34-35).'' The worst performers tended to newly-started,
medium-sized urban banks or NGOs in Africa or Asia (Heinen 2005: 39).

Without collateral fully backing many of the institution's financial
portfolios, this kind of risk becomes a major concern for prudent
investors. Recognizing moral hazard matters since, under moral hazard,
potential asset-poor customers typically face credit rationing---even
though they may have worthwhile investment projects. The projects are
worthwhile in the sense that they would generate enough income to allow
borrowers to in principle pay back loans with interest, and an
entrepreneur would profitably undertake the project with their own money
if they had the necessary fundsThe potential trouble is that the
borrowers cannot find funding for their projects. The same holds for
worthy institutions that are rationed out of the capital market due to
information problems. Much of the modern theory of financial
intermediation has been built upon this particular insight (Diamond
1984; Tirole, 2006) but the application of this theory to understanding
the role of public policy and social investors in shaping the emergent
financial intermediation landscape remains considerably unexplored.

We argue that commercial investors play important roles, but by virtue
of their philanthropic objectives, social investors can help increase
leverage, expand scale, and increase total output in a sector in ways
that commercial investors by themselves often cannot.

Specifically, we demonstrate cases where philanthropic-giving can in
principle raise output, increasing an economy's overall
productivity---and lead to GDP increases in addition to poverty
reduction. The possibility arises because social investors, by
definition, are willing to absorb costs in order to bring gains. This
is, of course, their reason for existence. They can do more than shift
resources: they can also help increase the total sum of resources
available. The question then is: How can they most powerfully enact such
``productivity-enhancing'' investments.\footnote{This language is from
  Pranab Bardhan, Samuel Bowles, and Herbert Gintis (2000).}

Our consideration of effective social investment begins with two ideas.
The first consideration is that philanthropic giving can expand the
contract space by being willing to subsidize the `limited liability
rents' that often stand in the way of more trades in the commercial loan
market. Limited liability rents refer to the returns that must be left
with borrowers, or the loan officers that monitor them so as to provide
incentives to effort when limited liability constraints restrict the
incentives that can be provided by threatening punishments for example
by using legal mechanisms to seize collateral. When such rents plus the
cost of loan monitoring become too large relative to expected project
gains, lenders may walk away from making loans to individuals without
collateral wealth, leaving socially valuable projects unfunded. A small
subsidy in such contexts can a times be enough to turn a small private
loans that might otherwise have been unprofitable now profitable. If
enough of those small loans become profitable, a new microfinance
institution might emerge where otherwise it might not have.

The second consideration is that effective giving may lead to multiplier
effects. It can ``crowd in'' other resources by catalyzing the
investments of others, increasing the recipient's ability to
``leverage'' its resources. Leverage (i.e. to be able to obtain funds to
use in addition to an institution's own equity) is critical as it
determines an institution's ability to reach wide scale, and wider scale
can lower average fixed costs. Effective social investment aims to
expand leverage, holding constant other social concerns.

In this framework, profitability is neither necessary nor sufficient for
attracting commercial capital. Merely earning profit will not generate
the leverage to take institutions to wide scale. Moreover, productivity
can be enhanced by supporting institutions that would not exist if
profitability was the sole concern.

The framework has a specific implication for the most effective use of
social investment. We describe arguments that push social investment
either ``high'' or ``low'' when directed to poor communities.
Specifically, we show that smart subsidy involves allocations that
either go to institutions serving relatively well-off (but still poor)
customers, where the potential for leverage is maximized, or to
institutions serving the poorest customers. There is less ``bang for the
buck'' in supporting a middle range of institutions that are
self-sufficient with their own resources but for which leverage (and
thus potential scale) is limited.

\hypertarget{the-model}{%
\section{The Model}\label{the-model}}

The analysis builds upon extensions to a simple model similar to what
Jean Tirole (2005) describes as a `workhorse model' for understanding
the modern theory of corporate finance and the role of net-worth in
shaping access to credit under moral hazard\footnote{Jean Tirole's
  sweeping 635-page \emph{Theory of Corporate Governance} textbook is
  built around a workhorse model very similar to this, which he calls
  the `simplest model of credit rationing ... to illustrate the role of
  net-worth (p. 114).' See also Holmstrom and Tirole (1997, 2008) and
  Conning (1999) for earlier analyses of the underlying model.}.Our aim
is to understand the range of financial contracts and intermediary
structures that can and will emerge to serve different segments of a
heterogeneous population of would-be borrowers, and how this financial
landscape might be transformed. To highlight the role of borrower net
worth we begin with the extreme assumption that would-be entrepreneurs
live in spatially segregated neighborhood according to their initial
asset holdings and that localized financial institutions (FIs) or bank
branches, managed by knowledgeable local managers, must be established
within each neighborhood in order to identify and monitor borrowers.
CCreditors to these FIs (which might include depositors) and, later in
the paper, informed equity investors may however make investments across
neighborhoods.

There will be up to four types of agents: (a) households that become
entrepreneurs to run small business firms; (b) one or more local
financial institutions in each neighborhood, wholly or partly owned and
managed by locally informed equity investors. These financial
institutions may lend out of their own equity capital and/or may be able
to intermediate funds from (c) private investors (possibly including
savings depositors). Finally, social investors may be able to affect the
nature and depth of the above relationsips through additional
investments, subsidies and guarantees of their own.

Each household is assumed risk neutral and to have access to a
production or trading project which requires lump sum investment
\emph{I} to be initiated. If the entrepreneur secures and invests the
lump sum and works diligently then the project succeeds and generates
cash returns \emph{X\textsubscript{i}} with probability \emph{p} or it
will fail and generate nothing with probability 1-\emph{p}. If this
borrower is non-diligent, for instance because she diverts effort or
resources from the project to other uses, then the probability of
success is reduced from \emph{p} to \emph{q} but the entrepreneur
obtains a \emph{private} benefit \emph{B}. Since the entrepreneurs
choice of diligence cannot be verified and private benefits cannot be
observed or seized by the lender, the potential for borrower moral
hazard arises

We assume potential entrepreneurs have no \emph{liquid} capital of their
own, so they must borrow the entire amount \(I\) to carry out the
project. Would-be borrowers may however have \emph{illiquid} assets
\(A\) that can serve as collateral guarantee in the event of failure.
These are assets the borrower prefers to keep tied up in other projects
or investments, or they may be pledgeable future cashflows. What matters
is that these resources can be credibly pledged to the lender in the
event that the financed project should fail. Although some microfinance
lenders do not explicitly require a collateral pledge there is little
doubt that most experienced loan officers form an estimate of how much
they believe can be credibly recovered from a borrower in the event of
low returns on the finance project, and this estimate of \(A_{i}\) will
be an essential determinant of the terms of the contracts
offered.\footnote{We assume here that the full value of may be taken
  from the borrower and costlessly transferred to the lender, but the
  analysis can be easily extended to cases where some or all of the
  value of is lost in the transfer.}

An MFI is assumed to begin with initial equity capital \emph{K} -- which
we shall often refer to as \emph{local intermediary capital}. The MFI
could be either a for-profit financial entity, in which case \emph{K}
would be interpreted as the owners' initial equity investment, or it
might be a not-for profit organization which is prohibited from
distributing dividends, in which case \emph{K} might be startup capital,
perhaps seeded by a donor.\footnote{In later sections we will expand our
  definition of to encompass certain types of outside equity or
  quasi-equity investments (such as subordinated debt) from donors or
  social investors.} \emph{K} is capital owned by, or under the control
of individuals who know how to operate an MFI in this particular
neighborhood and who will be in a good position to monitor the loans of
any borrowers that they may create. The point to emphasize is that \(K\)
is not easily created or simply borrowed from outside of the
neighborhood -- it is capital that is tied to people with a specific set
of skills. Much of the analysis that follows turns on understanding when
and how \(K\) can be leveraged to attract outside investors to expand
the scale of lending in a particular neighborhood.

If \(K\) is the MFI's only source of funds (i.e. no outside leverage is
possible), then after the MFI has paid for its fixed operating costs
\(F\), it will be able to lend out up to \(K - F\) worth of funds per
period and therefore reach \(N_{0} = (K - F)/I\) borrowers. Since the
MFI always has the option of keeping \(K\) invested in a safe investment
for gross return \(\gamma K,\) and it wants to remain sustainable into
the future we will assume that the MFI will only come into business if
after considering all cashflows (including possibly outside subsidies)
it can expect to earn this rate of return.

Any would-be entrepreneur has the option of not taking a loan and
engaging in an alternative activity or occupation where they are assumed
to earn a reservation payoff, normalized at zero. A financed project
yields either a verifiable return of \(X\)when it succeeds or it fails
to yield nothing. The probability distribution of project returns is
determined the borrowers' choice of effort or diligence. If the borrower
chooses to be diligent, in the sense that they exert personal effort and
use all borrowed funds for the financed activity as stated in the
business plan they described to their lender, then the project succeeds
with probability \(\text{p.}\) A borrower who is non-diligent lowers the
probability of success from \(p\) to \(q < p\), but captures some
private benefits \(B\) by avoiding effort and/or by redirecting some
borrowed funds to activities whose return may not observed or seized by
the lender. \(B\) therefore captures the opportunity cost of being
diligent.

A loan contract can be described as a contingent borrower reward
schedule \((s_{s},s_{f})\) where \(s_{i}\) indicates the net monetary
return to the borrower in state \(i\), which can be success (\(i = s\))
or failure (\(i = f\)).\footnote{With a few additional assumptions the
  model can be extended to a multiple-outcome multiple-action framework
  where the optimal contract takes the form of a standard debt contract.
  This complicates the analysis but does not fundamentally alter the
  main conclusions (Innes, 1990; Conning, 2006)} We assume that MFIs
offer competitive contracts that maximize the expected return to the
borrower. This assumption can be motivated either by assuming that each
MFI wants to maximize profits but that contestable entry pushes lenders'
profits to zero, and/or that MFIs have a social mission that leads them
to care about borrower welfare subject only to a sustainability (i.e.
break even) condition.

Consider a contract in which a borrower pledges all their available
collateral\(A\). We focus on this particular contract to identify the
`minimum collateral requirement' that a borrower would have to have in
order to be feasibly reached by a lender. The contract will take the
form \((s_{s},s_{f})\) =\((s, - A)\) where \(- A\) is the return to the
borrower when the project fails (i.e. the project generates zero returns
and the borrower loses the collateral pledge). The lender is repaid only
\(R_{f} = A\) when the project fails and \(R_{s} = X - s\) when it
succeeds. The lender will willingly participate in this contract so long
as the borrower is diligent and the expected return under diligence
\(E\lbrack R_{i}|p\rbrack = E\lbrack X_{i}|p\rbrack - E\lbrack s_{i}|p\rbrack\)
exceeds the cost of lending:

\(p(X - s) + (1 - p)A \geq \gamma I + F/N\)

where \(\gamma\) is the gross opportunity cost of funds to the MFI (one
plus the interest rate the lender pays) and \(F/N\) is the average fixed
cost per borrower. For the time-being we take \(\text{N\;}\)to be fixed,
but later endogenize its value. The line labeled \(PC(0)\) in figure 1
indicates combinations of \(s\) and \(A\) at which this lender's
participation constraint is exactly met. A point (\(- A,s\)) on this
Cartesian diagram represents a potential contract, but only contracts
that lie below \(PC(0)\) in this contract space satisfy the
participation constraint.

The expected return to the \emph{diligent} borrower is:

\[ps + (1 - p)( - A) - B\]

\(= p(s + A) - A - B\)

while the expected return to a \emph{non-diligent} borrower is:

\(q(s + A) - A\)

It follows that the borrower prefers to be diligent so long as (2)
exceeds (3) or

\(s \geq - A + \frac{B}{\Delta}\)

Where ∆\emph{=p-q}. This constraint restricts incentive-compatible
contracts to lie above line \(IC(0)\) in figure 1.\footnote{To keep
  things interesting we assume parameters such that the lender could not
  recover lending costs if the entrepreneur chose to be non-diligent.
  The lender will therefore only be interested in participating if they
  can be convinced that the contract satisfies this incentive
  constraint.} The range of feasible contracts must therefore lie in the
shaded region above \(IC(0)\) where the borrower remains diligent and
below \(PC(0)\) where the lender recovers all costs. Of particular
interest is the contract with the \emph{minimum collateral requirement},
or the lowest required pledge \(A\), since this contract will allow
reaching the largest group of borrowers\(.\) This contract is found at
the intersection of \(IC(0)\) and \(PC(0)\) and is labeled \(U\) in the
figure.

To see this mathematically note that limited liability sets a floor on
the borrower's return in the failure state of \(- A\). This and
constraint (4) requires that the borrowers' return to success not be set
below\(s = p/\Delta - A\) . A contract that satisfies limited liability
and incentive constraints therefore must leave the client with an
expected return of at least

\(E\left\lbrack s_{i}|p \right\rbrack = - A + \frac{\text{pB}}{\Delta}\)

Following Laffont and Rey (2003) we refer to expression \(pB/\Delta\) as
a \emph{limited liability rent} (LLR). It is the minimum share of
expected project returns that must be left to a client if incentives and
limited liability to be maintained. This rent is larger the larger is
the opportunity cost of diligence \(B\) and the more difficult it is to
infer diligence choice from simply observing the project outcome (i.e.
the larger is \emph{q/p})

Using (5) to substitute for \(E\left\lbrack s_{i}|p \right\rbrack\) in
(1) and then solving for the threshold value of \(A\) at which this
constraint exactly binds yields an expression for the \emph{minimum
collateral requirement} that a client would have to pledge to attract
the participation of an unsubsidized lender:

\((0,N) = \frac{\text{pB}}{\Delta} - \left\lbrack \text{pX} - \text{γI} \right\rbrack + \frac{F}{N}\)

Where \((0,N)\) refers to the minimum collateral requirement on a
zero-monitoring loan when the MFI services \emph{N} borrowers. Borrowers
with assets \(A \geq (0,N)\) will be able to obtain loans of size \(I\)
and earn expected return\(pX - \gamma I - F/N\), but entrepreneurs with
less than this amount will be excluded from the loan market even though
their projects may be socially profitable, i.e.
\(pX - \gamma I - F/N > 0.\) For later reference we define
\(= (0,\infty)\) to be the minimum collateral requirement for a
non-monitored loan in a really large MFI (i.e. where fixed costs per
borrower are essentially nil).

A higher level of \((0,N)\) implies a larger number of borrowers
excluded from the loan market. Expression (6) suggests the available
ways for MFIs and other financial innovators to improve loan access and
terms: (a) by acting to lower the size of the limited liability rent
either by acting to directly reduce the borrower's opportunity cost of
diligence \(B\) or changing the likelihood ratio (e.g. lending for
projects with lower \(l =\) \(q/p\) or that are `safer' under diligence
compared to non-diligence); (b) by lowering of average fixed costs
\(F/N\) either by lowering \(F\) or expanding scale by increasing the
number of borrowers \(N\); or by (c) raising the expected returns from
client's projects \(\text{pX}\) or lowering the MFI's cost of funds
\(\gamma\).

\hypertarget{monitored-lending}{%
\subsection{1.1. Monitored lending}\label{monitored-lending}}

A defining characteristic of most microcredit lending is the heavy use
monitoring and other incentive mechanisms to substitute for the use of
loan collateral guarantees. One way to interpret monitoring and control
activities is that they aim to try to raise \emph{p} relative to
\emph{q}, for instance via technical assistance that raises the
productivity of client efforts under diligence. This approach was
followed in Banerjee, Besley, and Guinnane (1994) and Madajewicz (2004).
An alternative interpretation is that monitoring of borrowers serves to
lower the value of \emph{B}. This is the approach of Holmstrom and
Tirole (1997), Conning (1999) and what we shall follow below.

`Monitoring and control activities' of this last kind aim to reduce the
scope for moral hazard by reducing the value of private benefits \(B\)
that the borrower stands to capture. A straightforward way to capture
the role of such active monitoring is to assume that by spending \(m\)
on monitoring and control activities a MFI can lower the borrowers'
benefits of non-diligence linearly from
\(B(0) = B\)to\(B(m) = B - \text{αm}\), where \(\alpha\) is a parameter
that captures the degree of monitoring effectiveness. We might think for
example of a lender making frequent and unannounced visits to an
entrepreneur's home or place of business, or their insistence that the
borrower prepare a business plan. These are just a few examples of the
costly monitoring activities that can be use to lower a borrower's scope
for diverting loan funds and other resources to uses other than the
financed activity. Since it doesn't matter whether we define \(B\) as a
private benefit to non-diligence or a cost to being diligent, we could
also think of monitoring as providing helpful technical assistance that
also lowers the borrower's cost of raising the probability of project
success.

Since monitoring activities are however themselves costly and difficult
activities, an MFI's financial backers must put in place the right
incentive and governance structures to make sure that their managers and
loan officers are properly motivated to carry out this monitoring at
optimal levels to pursue the objectives of its primary stakeholders. An
MFI's ability to leverage outside funds depends therefore in primarily
on how well they can convince outside investors that incentives and
monitoring effectiveness will not be damaged once the MFI begins to
scale up by using outside funds.

For the moment we shall set aside the question of outside funding and
focus first on MFIs that fund their operations using only the internal
equity funds \(K\) at their disposal. Monitoring by the MFI adds to the
total cost of lending (depicted as a shift in the participation
constraint from \(PC(0)\) to \(PC(m_{1})\) in figure 2 but it also
relaxes the incentive compatibility constraint (depicted as a shift from
\(IC(0)\) to \(IC(m_{1})\)). With monitoring the earlier expression (6)
for the minimum collateral requirement now becomes\footnote{For the time
  being, we continue to assume that \emph{N} remains fixed, but we add
  \emph{N} as an argument for later use.}:

\((m,N) = \frac{p(B - \alpha m)}{\Delta} - \left\lbrack \text{pX} - \text{γI} \right\rbrack + F/N + m\)

Under the assumption that on the margin monitoring is a worthwhile
activity or that \(\alpha/(1 - l) > 1,\) monitoring will lower the
minimum collateral requirement linearly. Figure 2 illustrates how an
increase of monitoring from \(m = 0\) to \(m = m_{1}\) might reduce the
minimum collateral requirement from \(= (0)\) to \((m_{1}).\) Additional
monitoring continues to reduce the minimum collateral requirement, but
more monitoring increases the cost of funds to the borrower who (from)
pays expected borrowing costs

\[\gamma I + F/N + m\]

It should be immediately clear that borrowers prefer less heavily
monitored loans because they are less costly. Hence borrowers will chose
loans with only as much monitoring as is just sufficient to bring the
minimum collateral requirement down to their available asset level\(A\).
Figure 3 depicts the relationship between monitoring and the minimum
collateral requirement . A borrower with available assets \(A < A\)
would choose a loan with just enough monitoring intensity \(m_{A}\) to
bring the collateral requirement down to their available assets, or
\((m_{A}) = A.\)

There is a limit to how far monitoring can expand access to asset-poor
borrowers. No borrower would voluntarily agree to a loan where the total
cost of borrowing exceeds expected project returns. This defines a
maximum monitoring level

\(= pX - \text{γI} - F/N\)

beyond which no borrower would willingly participate\(.\) To summarize,
monitoring can expand loan access to borrowers with assets in the range
\(A \in \lbrack(),(0)\rbrack\) that would have been excluded from the
non-monitored loan market.

\hypertarget{delegated-monitoring-leveraged-intermediation-and-scale}{%
\subsection{1.2. Delegated monitoring, leveraged intermediation and
scale}\label{delegated-monitoring-leveraged-intermediation-and-scale}}

Thus far we have said nothing of the organization or governance
structure of the MFI, where the MFI obtains its loanable funds from, or
what determines the scale of lending (i.e. the number or
borrowers\(N\)).

For each loan of size \(I\), amount \(I^{m}\) will be financed from MFI
capital (\(K - F\)) and amount \(I^{u} = I - I^{m}\) (possibly zero) by
raising funds from an uninformed outside investor (creditors and/or
depositors to the MFI). Raising outside funds offers the opportunity to
increase scale from \(N_{0} = (K - F)/I\) to \(N = (K - F)/I^{m}\)
borrowers.

Creditors or investors in the MFI might include local saving depositors
or national or international investors who make loans to the MFI. We
assume that investors have other opportunities for their capital, so
they will only be interested in investing if the MFI can credibly
guarantee an expected gross rate of return of \(\gamma\). The central
concern of any investor will be whether the management and loan officers
of the MFI have the proper incentive to diligently monitor borrowers so
as to preserve the quality of the loan portfolio.

There will now be three parties to a loan contract: the borrower, the
managers/owners of the MFI, and the outside investor(s). The timing of
the contracting problem is as follows: All parties agree to the terms of
the contract, the outside lender puts up \(I^{u}\)per borrower, the MFI
puts up \(I^{m} = I - I^{u}.\) The client then chooses his unobserved
diligence level, understanding and taking into account that, given their
stake in his project, the MFI will have incentive to monitor at
intensity \(m\) per borrower. Project returns are then realized and the
generated returns are then divided according to the terms of the
contract.

Loan monitoring is itself an activity that is subject to moral hazard.
Since the MFI management and staff's actions cannot in general be
directly verified by outside investors incentives to monitor must be
provided by tying the MFI agents' rewards to observable outcomes on the
borrowers' projects. The way outside investors will be convinced that
the MFI owners and their staff have sufficient incentives to maintain
the quality of the loan portfolio is by demanding that the MFI place
enough of their own capital \(I^{m}\) at risk in each borrower's project
so as to have incentive to monitor by enough to keep the borrowers
diligent.

The MFI's will earn net return \(R^{m} - \gamma I^{m} - m\) if it
monitors and the project succeeds and \(- \gamma I^{m} - m\) if it
fails.\footnote{For the moment we focus on the terms of the contract as
  if there were only one borrower. In practice the lender will not want
  to tie the MFI's return to the outcome of a particular project, but
  rather to the outcomes from the entire portfolio. We will return to
  this desire for diversification later. For the moment, and to
  simplify, we follow Holmstrom and Tirole (1997) in assuming that all
  of the borrowers have projects with stochastic but perfectly
  correlated outcomes, so all projects succeed or fail at the same time.
  We will later return to the more general case with diversification, as
  studied by Conning (2005), Tirole (2005, p158-163) and Diamond (1984).}
To have incentives to monitor the MFI must expect to earn more by
monitoring at intensity \(m\) than by not monitoring at all, where
intensity \(m\) is the minimum amount required to induce the borrower to
be diligent, given the other terms of the contract. The MFI's incentive
compatibility constraint is therefore:

\[pR^{m} - \gamma I^{m} - m \geq qR^{m} - \gamma I^{m}\]

\(R^{m} \geq m/\Delta\)

The MFI must also be willing to participate, which is guaranteed if they
can expect to earn as much from their efforts and funds inside the
relationship as outside (i.e. they expect to earn as much by investing
\(I^{m}\) elsewhere as they do by lending:

\(pR^{m} \geq \gamma I^{m} + m + F/N\)

Since the MFI is responsible for organizing the banking activity,
including renting buildings and equipment, they must cover \emph{F/N} to
enter and stay in the financial intermediation business.

For the moment we will assume that MFIs are driven by a social concern
for the welfare of the borrower (or equivalently that competition
amongst potential intermediaries) so there are no economic profits or,
equivalently the return on MFI equity is equal to\(\gamma\). We will
later relax this assumption, but for the moment the assumption of zero
profits implies that offered loan contract will be on the most favorable
terms the borrower, so this participation constraint as well as that of
the outside investor will bind in equilibrium. Substituting the
intermediary's binding incentive constraint in this participation
constraint yields:

\(\gamma I^{m}(m) = \frac{\text{qm}}{\Delta} - \frac{F}{N}\)

which defines the MFI's required equity stake in the entrepreneur's
project if incentives are to be maintained to monitor at intensity
\emph{m} and lending costs are to be minimized to maximize leverage and
outreach.\footnote{We later explore the possibility that intermediary
  capital is scarce and MFIs earn a return on equity above γ.} The
higher the required amount of monitoring \emph{m} to keep a borrower
diligent, the higher the stake the MFI must hold in the project since
\emph{q}/∆\textgreater{}0. This expression suggests an immediate
explanation for why leverage has remained disappointingly low in
microfinance. As Ledgerwood (1998, page 41) explains `most MFIs are not
fully leveraged due to their inability to borrow funds on the basis of
their performance and debt capacity.' The above expression suggests why:
MFIs must monitor their loan portfolios heavily to maintain its value.
Monitoring is a replacement for missing collateral. The more heavily the
MFI concentrates on a segment of the market without collateralized loans
(which could be sold to debt collectors) the more the portfolio value
depends on the level of monitoring, and monitoring incentives require
the MFI retain a deeper stake and therefore that it be less leveraged.
The commonly held myth that the pursuit of higher MFI profits should
translate into higher leverage does not take this factor into account.

In a situation where MFIs seek to maximize borrower returns (or in a
competitive market) both the MFI intermediary and the uninformed outside
lenders (if any can be attracted) earn just enough to cover their costs
and earn γ in expectation on their investments. The minimum collateral
requirement can be found as before, by substituting the borrower's
binding incentive constraints the MFI's monitoring incentive constraint
and the uninformed lender's binding participation constraint
(\(pR_{u} + (1 - p)A = \gamma I_{u}\)) into the MFIs participation
constraint:

\(pR_{m} \geq \gamma I_{m} + m + F/N
pX - \left\lbrack \frac{p(B_{0} - \alpha m)}{\Delta} - A \right\rbrack - \gamma I_{u} \geq \gamma I_{m} + m + F/N\)

and solving for \emph{A} the collateral minimizing solution is again as
in (7):

\((m,N) = \frac{p(B - \alpha m)}{\Delta} - \left\lbrack \text{pX} - \text{γI} \right\rbrack + F/N + m\)

with the important new interpretation that \emph{N} will now be treated
as an endogenous variable. Higher leverage will mean the MFI can reach a
larger number of borrowers \(N\)with any given level of equity
\(K - F\). This will reduce the average fixed cost of borrowing per
borrower \(F/N\) which will, in turn, open up space to adjust the terms
of the financial contracts and reach yet more borrowers.

For borrowers in asset class \emph{A} we can now solve for the
equilibrium level of borrowers \emph{N} and the minimum required amount
of monitoring per borrower \emph{m}. These will be given by the solution
to the following system of equations:

\((m,N) = A\)

\(N = \frac{(K - F)}{I^{m}}\)

where \(I^{m}\) is given by (11). As demonstrated in the appendix this
system can be solved to give the following closed-form solution for
\emph{N} and \emph{m} as a function of a client's initial asset level
\(A\) and parameters describing the MFI's lending technology and equity
base:

\(N^{*} = \frac{\Gamma}{\left\lbrack - A \right\rbrack}\mspace{6mu}\)

\(\mspace{6mu} m^{*} = \left\lbrack - A \right\rbrack \cdot \frac{\Delta}{q}\frac{(\gamma K - (1 - \gamma)F)}{\Gamma}\)

where \(= (0,\infty)\) and
\(\Gamma = \frac{((1 - \alpha)p - q)}{q}\lbrack\gamma K - (1 - \gamma)F\rbrack - F\).

Required monitoring falls towards zero as the client's collateral wealth
\emph{A} approaches\(\), the collateral requirement for an non-monitored
loan in a big bank (where fixed costs per loan are nil). Leverage, and
therefore also the number of borrowers \emph{N} that can be reached,
rise exponentially as \emph{A} approaches\(\)from below. This apparently
extreme result -- that leverage becomes infinite -- is an artifact of
our implicit assumption that it is costless for the MFI or any third
party to collect collateralized debt. In practice local MFIs will have
an advantage in collecting compared to an outside investor and there
will be moral hazard involved in the collection of such debt if
collection is at all costly (Townsend, 1979; Diamond, 1984). It is also
true that uninformed lenders (small depositors in particular) will
typically not be in a good position to observe the MFI's debt to equity
holdings. For all these reasons financial institutions that accept
deposits are typically regulated which in effect places an institutional
cap on maximum leverage. To capture this idea and to keep the number of
borrowers from rising to infinity we will impose an institutional limit
on leverage in the form of \(N \leq N_{\max}\) This is without loss of
generality because the focus of our analysis is on MFIs that will be be
unable to reach this high level of leverage anyway.

Figure 4 shows how the equilibrium level of \emph{N} and \emph{m} would
be determined for MFIs that aim to reach target groups with assets
\emph{A} and \emph{A'}.

The MFI's balance sheet can represented simply as:

\begin{longtable}[]{@{}ll@{}}
\toprule
\textbf{Assets} & \textbf{Liabilities}\tabularnewline
\midrule
\endhead
\begin{minipage}[t]{0.47\columnwidth}\raggedright
\(N \cdot I\mspace{6mu}\)Loans\strut
\end{minipage} & \begin{minipage}[t]{0.47\columnwidth}\raggedright
\(N \cdot I^{u}\mspace{6mu}\) Debt

\(N \cdot I^{m} = K - F\) Owners Equity\strut
\end{minipage}\tabularnewline
\bottomrule
\end{longtable}

Of course, when the MFI is unable to lever up and attract outside
deposits or other forms of debt the liability side of the balance sheet
will consist solely of owners equity (including donations) and the
number of loans will simply be \(N_{0} = (K - F)/I\).

A numeric example will help to illustrate the relationships captured by
the model and predictions regarding the reach of different types of
lender on this market. Consider a simple case where the parameters are
as given in Table 1.

Consider first the earlier case where the MFI lends only out of its
available capital base \(K,\) without leverage, so \(N = N_{0}\). For
these parameters the moral hazard problem would lead lenders to demand
\(= (0;N_{0}) = 130\) units of pledgeable income to cover a
zero-monitoring loan of \(I = 100\). This is the kind of lending carried
out by ordinary banks that typically work only with clients with
substantial savings or cash flows relative to the size of the borrowed
amount. The associated loan would specify fixed repayment amounts of
\(= (130,130)\) on a loan of size \(I = 100.\)

The implied expected cost of funds to the borrower is forty percent:
twenty percent to cover the lender's variable cost of funds
\(\text{γI}\) and another \(10\) percent to cover average fixed cost of
\(F/N_{0} = 10\) per loan. The borrower earns an expected return
\(pX - \gamma I - F/N_{0},\) which for these parameters is
\(30( = 180 - 120 - 30).\)

Consider the possibilities and limits of monitored lending to reach more
borrowers. Any borrower with less pledgeable assets \(A\) will have to
resort to a more expensive monitored loan, or receive no loan at all.
Since monitoring raises the cost of borrowing, borrowers will want a
loan with only just as much monitoring as is minimally required to bring
the collateral requirement down in line with their available assets
\(A\), or the minimum \(m\) to set \((m,N_{0}) = A\) using (7). For
instance a borrower with \(A = 60\) would require a loan with minimum
monitoring intensity \(m = 35,\) assuming an MFI with no leverage, or
\(N = N_{0}\).The borrowers total cost of borrowing \(I\) will then be
\(\gamma I + F/N_{0} + m = 120 + 20 + 35 = 175\), for an implied
interest rate of 75 percent.\footnote{While this borrower pays out only
  60 out of collateral when their project fails (compared to 130 for the
  earlier borrower who could post collateral to receive a loan with zero
  monitoring), they take home only 23.33 when the project succeeds
  (compared to 70). This leaves this borrower with an expected profit of
  only 15 (compared to 30).}

Hence, although monitoring expands access, these asset-poor borrowers
pay much higher financing costs to cover the MFI's higher costs of
monitoring. For the parameters above, the maximal feasible level of
monitoring is \(= 50\) and the lowest the collateral requirement can be
pushed is down to \(()\) = 30. Monitoring carried out by the MFI with no
leverage can therefore expand access to borrowers with assets in the
range \(A \in \lbrack\)30\(,\) 130\(\rbrack\) who would otherwise not
have been served. Any borrower with \(A\) less than \(()\) will prefer
to forgo a loan and remain in the low-return occupation. Even though
they have access to a socially profitable investment (and one that they
would carry out if they had had their own funds) they are unable to
secure financing to profitably carry out the project due to the the
nature of the information asymmetry.

In the example thus far we have assumed that the MFI uses only its own
funds. It has \(K = 11000\) of capital, \(F = 1000\) of annual fixed
costs, and it leverages no outside funds. It was therefore able to lend
out only 10,000 through \(N_{0} = 100\) loans at \(I^{m} =\)100 dollars
per loan). This MFI's balance sheet would simply look as follows:

\begin{longtable}[]{@{}ll@{}}
\toprule
\textbf{Assets} & \textbf{Liabilities}\tabularnewline
\midrule
\endhead
10,000 (=\(N_{0}I^{m})\mspace{6mu}\) Loans & 10,000 \(( = K - F)\)
Owners Equity\tabularnewline
\bottomrule
\end{longtable}

Consider next the case of an MFI that attempts to leverage outsider's
funds to increase outreach and lower average fixed costs. Notice
immediately from (11) that if there is no monitoring required there is
no minimum required stake \(I^{m}\) and leverage \((I - I^{m})/I^{m}\)
could in principle be infinite. In that case the MFI would be a perfect
and costless intermediary who could always repay its investors, because
it has fully collateralized loans and zero costs of intermediation. This
would be the financial equivalent of a perpetual motion machine. Since,
as previously discussed, even fully collateralized loans require some
monitoring and collection is itself a costly activity, in reality
\(I^{m}\) would never be zero, if for no other reason than that
regulators would intervene to limit leverage via capital adequacy
requirements. We will assume that regulators require MFIs to have a
least ten cents of their own capital at risk for every dollar loaned
(i.e. \(I^{m} \geq 0.1I\)) \(,\) or that they can borrow no more than
ten times their equity base.

For the above parameterized example this means that \(N_{\max}\) With
this many borrowers the minimum collateral pledge on a non-monitored
loan drops to \(= (0;N_{\max}\) (from 130 when\(N = N_{0} = 100\)) on
account of the fact that the average fixed cost per borrower has been
reduced from \(F/N_{0} = 10\) to \(F/N_{\max}\).

A collateral requirement of 121 on a loan of size 100 still leaves a lot
of excluded borrowers, so let's turn to how increased leverage expands
the possibilities of using monitored loans. The optimal monitoring
intensity \(m\) and maximum feasible scale \(N\) is determined by
equations (16) and (15) and will depend on the level of \(A\) of the
target borrower group. Starting from the zero-monitoring collateral
requirement \(,\) we can lower the collateral requirement to reach
successively poorer borrowers with assets \(A <\) by increasing
monitoring intensity \(\text{m.}\) As monitoring intensity rises
however, the minimum investment stake of the MFI \(I^{m}\) must also
rise (given by (11)). This means that MFIs that specialize in reaching
neighborhoods with lower \(A\) achieve lower leverage \(I/I^{m}\) and
therefore also that they reach a lower number of borrowers \(N\) with a
given amount of startup capital (expression ((14).

Figure 4 describes the financial landscape that emerges in this model.
We assume that there are many borrowers at each asset level in each
neighborhood and a different MFI (with initial capital \(K\) and fixed
costs \emph{F}) specialized in lending to borrowers in each of these
neighborhoods. In the emergent landscape neighborhoods will be
partitioned into four groups, depending on the initial asset level A of
the household in each neighborhood. The figures show the relationships
between pledgeable assets \(A\) of the MFI target group borrowers and A)
the number of borrowers that the MFI can reach with its available
capital \(K,\) B) net returns per borrower, and C) total returns (net
returns per borrower times number of borrowers). Neighborhoods will be
characterized as falling into one of the following groups:

\emph{\textbf{{Group I} -- Excluded Borrowers}} (\(A < A_{\text{LO}}\))

Entrepreneurs with assets \(A\) below threshold level
\(A_{\text{LO}} = (,N_{0})\) cannot obtain financing at any profitable
rate. As a rule MFIs can only reach asset-poor borrowers via increased
monitoring, but these borrowers are so poor that the monitoring required
would be prohibitively expensive. The threshold level of monitoring was
earlier described in (8), as \(= pX - \gamma I - F/N_{0}.\) In the
absence of any subsidies MFIs that try to reach borrowers with \(A <\)
\(A_{\text{LO}}\) would lose money and eventually go out of business.

\emph{\textbf{{Group II} -- Monitored loans from non-leveraged MFIs}}
(\(A \in \lbrack A_{\text{LO}},A_{\text{LEV}}\rbrack\))

MFIs that service borrowers with assets in this range use monitored
lending but cannot attract any outside investors because the minimum
required stake \(I^{m}\) in (11) that is required to attract an outside
investor exceeds the total value of the loan \(I\). These MFIs can reach
asset-poor borrowers using their own capital, or \(I^{m} = I\) but these
MFIs can reach only as many borrowers \(N_{0} = (K - F)/I\) as their
capital permits. The net expected return to a borrower is rising with
asset levels over this range (since borrowers with more collateral
require less expensive monitoring), as depicted in the top panel of the
figure.

These MFIs cannot sell any part of their portfolio of monitored loans
without, in the eyes of outside investors, diluting their incentives to
monitor the quality of that loan portfolio diligently. MFIs trying to
reach borrowers with assets in this range cannot attract outside
investors not because they are not profitable (in fact they are earning
normal profits, as are all other MFIs). The commonly heard claim that
leverage would be raised if only the MFIs raise profits further by
raising the borrowers' required repayments can be seen to be flatly
untrue here. In fact, demanding a higher repayment of borrowers would
only exacerbate rather than improve the situation, by aggravating the
moral hazard problem for borrowers.

\emph{\textbf{{Group III} -- Monitored Loans from leveraged MFIs}}
(\(A \in \lbrack A_{\text{LEV}},A_{\text{HI}}\rbrack\))

MFIs that service borrowers in this range are able to leverage their
capital, and therefore reach a larger number of borrowers. This is
because the required minimum equity stake \(I^{m}\) is less than the
full amount of the loan \(\text{I.}\) These MFIs can therefore reach
\(N = (K - F)/I^{m} > N_{0}\) borrowers, and this number of borrowers
increases exponentially with higher the level of \(A\) (up to
\(= (0;N_{\max}\) as depicted in the middle panel of the figure. The
borrower's net return rises more quickly with \(A\) than for borrowers
served by non-leveraged MFIs (top panel of the figure) because the costs
of borrowing fall more rapidly with rising \(A\) because these MFIs can
take advantage of increasing scale (which lowers average fixed costs).
The cutoff level \(A_{\text{LEV}} = (m_{\text{LEV}};N_{0})\) where
\(m_{\text{LEV}}\) is given by the \(m\) at which
\(\gamma I = \frac{\text{pm}}{(p - q)} - m\).

\emph{\textbf{{Group IV} -- Non-Monitored Loans from highly leveraged
MFIs}} (\(A \geq A_{\text{HI}}\))

The lower panel of figure 4 shows total borrower gains from the presence
of one MFI with \(K = 1000\) in each neighborhood. More value is created
in the less asset-poor communities because MFIs that operate there have
lower costs due to less monitoring and larger scale due to the
possibility of outside leverage. This is what the financial structure of
many markets looks like: the private financial system serves lots of
people in those segments of the market where borrowers have lots of
assets or plegeable income (e.g. stable wages that can be garnished) but
fewer or no borrowers as we move `downstream' to segments of the market
where borrowers have less pledgeable assets, because they are asset poor
or self-employed. Of course we might not expect there to be neccessarily
the same amount of `intermediary capital' \(K\) available in each
segment of the market, but we will start with that assumption in order
to explore how additional contributions of capital or subsidy will have
an impact on different segments of the market.

\hypertarget{smart-subsidies}{%
\subsection{1.3. `Smart' subsidies}\label{smart-subsidies}}

The graphs point to a major challenge confronting microfinance
institutions and their financial backers. Donors may be inclined to
target capital toward MFIs that focus upon asset-poor clients and the
poor because donors typically care more about the social welfare of poor
individuals, but the figure and the model suggests that each extra
dollar of capital transferred to higher-end MFIs can `crowd in' more
private sector resources than the same dollar targeted to an MFI trying
to reach the poor.

Before examining how a philanthropist or donor might choose when faced
with these tradeoffs and whether it might choose to direct its resources
via subsidies, grants or loans, it will be worthwhile to briefly pause
to clarify a few terms. There is a great deal of debate and sometimes
confusion in the policy literature over the role of subsidies. History,
theory and evidence indicate clearly that subsidies often do more harm
than good. A badly designed subsidy leads to a soft budget constraint
problem which creates dependence and weakens rather than strengthens MFI
incentives to look after loan portfolio quality. But in models such as
this where problems of asymmetric information and limited liability
create market failures, a carefully chosen subsidy can create markets
where they did not previously exist and crowd in resources in ways that
may raise rather than lower total surplus.

There may therefore be `smart' ways to deliver subsidies and a primary
aim of this framework is to help identify where such `smart subsidies'
might be possible. Consider for example how optimal contracts could be
adjusted if an MFI stepped in to subsidize fixed costs. So long as the
subsidy is delivered in a way that leads parties to recontract to
maintain incentives then outside investment, MFI outreach and borrower
welfare could well be enhanced. To fix ideas, consider the case of a
donor that had \$500 funds to partially subsidize the fixed costs of
just one MFI, reducing these from \(F = 1000\) to \(F^{'} = 500.\) Which
MFI should the philanthropist target for `maximum impact?'

In debating this issue it is important to draw a distinction between
levels and changes. The panels in figure 4 indicate levels and clearly
show that (in the absence of subsidies) an MFI with \$1000 of capital
\(K\) focusing on asset-rich borrowers can reach more borrowers
\((N_{A})\) and generate more total net surplus
(\(G_{A} = pX - \gamma I - m_{A} - F/N_{A})\) than another MFI with the
same \$1000 of capital \(K\) focusing on asset-poor borrowers.

This might seem to suggest that if a philanthropist had another \$500 to
allocate, they might want to deliver the subsidy to an MFI that targets
the relatively more affluent, since each dollar of subsidy can be used
to `leverage' outside funds. But this ignores the fact that the \$500
can also `leverage' capital in a different way. Suppose that in a poor
community there are local investors with local intermediary capital
\(K\) who, given existing fixed costs, information problems and the
average net worth of borrowers in this neighborhood, have not found it
profitable to pool these resources to form an MFI, but that with a \$500
subsidy to fixed costs these investors could earn enough to want to
enter business. In this case the \$500 subsidy would have 'leveraged' in
local capital that brought into existence an MFI that can serve
\(N_{0}\) borrowers, where no borrowers were previously attended.

The panels in Figures 5 and 6 illustrate the possibilities using our
parameterized example. The panels in figure 5 reproduce those in figure
4 (lines in red), but now also draw in blue dotted lines that represent
outcomes from an alternative scenario where each of the MFIs (or
potential MFIs arranged according to the asset levels \(A\) of the
borrowers they attend) would respond to a subsidy that lowered its fixed
costs from 1000 to 500. In analyzing the effects of these interventions
it will be useful to classify MFIs and the neighborhoods \(A\) in which
they reside into `impact groups' according to the type of transformation
they undergo as a consequence of an intervention:

\(\begin{matrix}
\mathrm{\text{Impact}}\mspace{6mu}\mathrm{\text{Group}} & \mathrm{\text{Definition}} & \mathrm{\text{Description}} \\
A & \mathrm{\text{Group}}\mspace{6mu} I\mspace{6mu} \rightarrow \mathrm{\text{Group}}\mspace{6mu}\mathrm{\text{II}} & \begin{matrix}
\mathrm{\text{New}}\mspace{6mu}\mathrm{\text{MFIs}}\mspace{6mu}\mathrm{\text{emerge}}\mspace{6mu}\mathrm{\text{where}}\mspace{6mu} \\
\mathrm{\text{none}}\mspace{6mu}\mathrm{\text{previously}}\mspace{6mu}\mathrm{\text{existed}}.\mspace{6mu} \\
\end{matrix} \\
B & \mathrm{\text{Group}}\mspace{6mu}\mathrm{\text{II}}\mspace{6mu} \rightarrow \mathrm{\text{Group}}\mspace{6mu}\mathrm{\text{II}} & \begin{matrix}
\mathrm{\text{Supports}}\mspace{6mu}\mathrm{\text{existing}}\mspace{6mu}\mathrm{\text{MFIs}},\mspace{6mu} \\
\mathrm{\text{no}}\mspace{6mu}\mathrm{\text{new}}\mspace{6mu}\mathrm{\text{leverage}}. \\
\end{matrix} \\
C & \mathrm{\text{Group}}\mspace{6mu}\mathrm{\text{II}}\mspace{6mu} \rightarrow \mathrm{\text{Group}}\mspace{6mu}\mathrm{\text{III}} & \begin{matrix}
\mathrm{\text{Transforms}}\mspace{6mu}\mathrm{\text{non}} - \mathrm{\text{leveraged}}\mspace{6mu} \\
\mathrm{\text{into}}\mspace{6mu}\mathrm{\text{leveraged}}\mspace{6mu}\mathrm{\text{MFI}}. \\
\end{matrix} \\
D & \mathrm{\text{Group}}\mspace{6mu}\mathrm{\text{III}} \rightarrow \mathrm{\text{Group}}\mspace{6mu}\mathrm{\text{III}}\mspace{6mu}\mathrm{\text{or}}\mspace{6mu}\mathrm{\text{IV}} & \begin{matrix}
\mathrm{\text{Increased}}\mspace{6mu}\mathrm{\text{leverage}}\mspace{6mu}\mathrm{\text{in}}\mspace{6mu} \\
\mathrm{\text{already}}\mspace{6mu}\mathrm{\text{leveraged}}\mspace{6mu}\mathrm{\text{MFI}}. \\
\end{matrix} \\
\end{matrix}\)

As can be seen a subsidy delivered to every neighborhood would shifts
\(A_{\text{LO}}\) to the left, with the effect of creating impact group
A by bringing into existence new MFIs to serve asset-poor neighborhoods
where none existed before. It also lowers \(A_{\text{LEV}}\), allowing
previously non-leveraged MFIs in (Impact Group C) to now access the
private debt markets for the first time. Through subsidies (in the form
or outright grants, subsidized loans, loan guarantees, etc) donors can
bring new MFIs into existence, they can help transform existing
non-leveraged MFIs into leveraged intermediaries that are now better
able to raise funds from private markets. Finally they can help increase
leverage in MFIs that can already raise private debt.

The question remains how should donors decide to target their scarce
funds across differnt MFIs and Impact groups? The incremental impact of
these subsidies can be seen in figures 5 and 6. Figure five shows
variable levels before and after the intervention, while figure six
looks at this difference, and decomposes it into component parts. The
top panel of figure 6 shows \(\Delta G_{A}\), the marginal impact of the
subsidy on the net returns of a single borrower in each neighborhood,
where \(G_{A} = pX - \gamma I - m_{A} - F/N_{A}\) and
\(\Delta G_{A} = G_{A}^{'} - G_{A}.\) The next panel shows the number of
new clients created in each neighborhood \((N_{A}^{'} - N_{A})\) and the
third panel shows the total gain from the intervention
\(G_{A}^{'} \cdot N_{A}^{'} - G_{A} \cdot N_{A}\), in each segment. This
can in turn be decomposed into gains to existing clients, given by the
number of existing borrowers in that segment \(N_{A}\) times the change
in net return (\(G_{A}^{'} - G_{A}\)) per borrower, and gains to new
clients, given by the change in number of clients
(\(N_{A}^{'} - N_{A}\)) times the net return to each of these new
clients \(G_{A}^{'}\) or:

\[{G_{A}^{'} \cdot N_{A}^{'} - G_{A} \cdot N_{A}
}{= N_{A} \cdot (G_{A}^{'} - G_{A}) + G_{A}^{'} \cdot (N_{A}^{'} - N_{A})}\]

The panels shows that the impact of a subsidy on a representative
borrower in each MFI is not at all uniform across neighborhoods and
impact groups. Even though the new MFIs created in the lower Group A
segment of the market are not able to leverage funds from outside
sources, and they could not survive without the subsidy, the donor's
funds have `crowded-in' \emph{local} intermediary capital, and by doing
so have helped borrowers with socially profitable projects to become
entrepreneurs. As can be seen many more new clients were created by
targeting this subsidy to Impact Group A neighborhoods than to Impact
Group B and C and many (but not all) Group D neighborhoods. If a donor's
main objective is to brag about the number of new clients its
intervention has helped to create then donors should focus their efforts
on MFIs in poor neighborhoods. The number of new clients in
neighborhoods in Group B hardly increases at all, because these MFIs do
not achieve leveraged. The subsidy allows the MFI to increase the pool
of funds available for lending by 500, allowing five new borrowers to be
reached with a \(I = 100\) loan each. The impact is larger in all other
cases because the intervention leverages in additional sources of
financing to expand scale.

Any proper analysis of impacts should however not just count new
borrowers but also consider the impact of the intervention on individual
borrowers' incomes. As the bottom panel of figure six indicates the per
borrower impact group A rises rapidly with the borrowers' initial asset
holdings\(A\). This is because the poorest amongst this group of all new
borrowers will have most of their expected project returns eaten up by
the very high monitoring costs needed to bring them into the market. The
return left to the borrower therefore rises as monitoring costs drop
linearly with higher \(A\).

Focusing on the next impact group B (i.e. borrowers served by MFIs that
began as un-leveraged Group II neighborhoods and remain unleveraged
after the intervention) we find by looking at the lower panel that most
of the gains come from extending a small gain to a an existing group of
borrowers, rather than from new borrowers. That is because these are
borrowers who already had access to capital, and since no leverage is
achieved, scale is increased only very modestly. Each of these borrowers
is benefited by the fact that the average fixed cost per loan is reduced
by \(500/N_{0}( = 5)\) per borrower plus an additional amount of savings
from the need for a bit less monitoring at each level of \(A\). In
Impact Group C on the other hand the number of borrowers starts to rise
more rapidly with \(A\) because the multiplicative impact of leverage.
As drawn, and for these parameters, Group C is not a very large number
of borrowers.

The analysis of group D is interesting. This impact group consists of
MFI that already achieved leverage, but with the aid of a subsidy are
now able leverage further. The multiplier effect of the subsidy via
leverage shows up primarily as an exponentially rising increase in the
number of new clients. The impact per client is however declining with
\(\text{A.}\) This is because were already in very highly leveraged
institutions where average fixed costs per borrower were low and
borrowers were already not very heavily monitored. Increased leverage
therefore shaves relatively little off the already low cost of reaching
these borrowers. As drawn however this dwindling impact is being
multiplied by a fast rising population of new borrowers to yield an
overall rising total impact of the subsidy as one targets more affluent
neighborhoods.

The graph on the lower panel of figure 6 also draws in a horizontal line
at 500. If we were being strict utilitarians and measuring benefits as
unweighted dollar value of total gains by neighborhood, we would
conclude that expected benefits to society of the intervention do not
compensate for the 500 subsidy cost for the very poorest borrowers. Of
course this will not typically be the welfare criterion a donor might
choose to pick since a donor is more likely to place higher weight on
gains created in poorer households.

Note that the free market outcome in the absence of philanthropists and
their subsidies (including subsidized investments) is constrained
efficient. That is, given the initial distribution of assets, structure
of the economy including information asymmetries, and the assumption of
voluntary exchange, there are no pareto-improving redistributions.
Entrepreneurs cannot commit to repaying others in society for loans or
transfers that can raise total output, because they cannot commit to
making such repayments without undermining their own incentives to
maintain the required level of effort. When however a new Bill Gates
emerges to become a philanthropist or a donor however, total output can
be increased because he does not demand full repayment (and indeed the
transfer would be Pareto-improving since the subsidized borrower gains
and Bill Gates is made no worse off since the contribution was
voluntary).

\hypertarget{profits-ownership-and-incentives}{%
\subsection{1.4. Profits, ownership and incentives
}\label{profits-ownership-and-incentives}}

\hypertarget{this-section-very-rough}{%
\subsection{(this section very rough)}\label{this-section-very-rough}}

So far we have assumed the MFI behaves so as to maximize outreach and
with that lower costs and increase benefits to pass on all the gains in
surplus to the average borrower. This would be consistent with the
behavior of a non-profit that cannot distribute profits and whose
incentives are aligned to maximize borrower welfare because of the
`social mission' orientation of its principal shareholders. The MFI is
able to sustain itself over time but the rate of return on equity is γ,
the same rate return as investments in other areas of the economy.

An intermediary with a profit motive will seek to earn an expected rate
of return on equity in excess of γ.

\(p\frac{m}{\Delta} = \beta I_{m} + m + \frac{F}{N}\)

We can interpret β as the expected return on equity:

The ROE is γ when intermediary's IC constraint binds. If however
intermediary capital is scarce, they will be able to earn a rent
β\textgreater{}γ. This will be associated with a lower
\emph{I\textsubscript{m }}(i.e. higher leverage). Appendix 2 explores
this scenario in further detail\ldots{}

\hypertarget{extensions}{%
\subsection{1.5. Extensions}\label{extensions}}

As in all models we have made simplifying assumptions to highlight
important mechanisms, but these simplifications may at times bias the
conclusions. A case in point is the assumption we have made that within
each neighborhood resides an unlimited pool of identical entrepreneurs
each with a high social return project (i.e. one that generates 200 for
success with probability 0.9 when the borrower is diligent). This
assumption tilted the conclusion of our analysis to favor subsidies and
social investment toward higher-leveraged MFIs because under these
assumptions these MFIs could pull in a large number of new borrowers. A
more realistic assumption is that the distribution of returns across
borrowers is heterogeneous, and that the marginal quality of the
entrepreneur's project declines as the MFI expands scale in a particular
neighborhood.

Under this new set of assumptions the social investors' tradeoffs are
likely to be rather different. Suppose for example that instead of all
borrowers having a project with a certain return of \(X = 200\) under
`success' this return where itself a random number drawn from a
distribution. This distribution is assumed to be identical across
neighborhoods and lenders in each neighborhood can observe \(X\) prior
to giving a loan. It is easy to see that under the private market
solution with no subsidies, all else equal, clients with lower \(X\)
will receive more heavily monitored loans, and lenders will extend scale
by moving from higher \(X\) borrowers to lower \(X\) borrowers to
maximize total surplus per neighborhood. It should also be clear
therefore that in a higher \(A\)neighborhood the last borrower to be
reached has a lower \(X\) compared to a poorer neighborhood where there
is less leverage. Under these circumstances the tradeoffs faced by a
donor considering where to place a subsidy are now different. They are
now more likely to tilt their decision toward poorer neighborhoods
because that is where they will find borrowers with higher expected
social marginal returns.

\hypertarget{conclusion}{%
\section{2. Conclusion}\label{conclusion}}

The act of making an investment in a microfinance institution typically
changes the incentives for the institution's owners and managers. The
investors now take on risk that had previously been borne more heavily
by the microfinance institutions. The institutions and their managers
now face fewer (or different) consequences of their actions. This shift
in who bears the risk for successes and failures can seriously diminish
incentives for diligence on the part of institutions and their
staff---such that a past history of successes may be a poor predictor of
future performance.

Drawing on the discussion of moral hazard above, we created a landscape
of outcomes. In a free market system, without social investors, we
expect that the landscape of finance for the poor would be divided into
four broad situations.

1.  \emph{The highly-leveraged.} The first involves institutions lending
to better-off borrowers that have collateral---and where courts function
effectively to seize assets. Here, microfinance should be maximally
leveraged since outside investors can invest with limited concern about
the diligence of the institution's staff. Loans do not have to be
monitored in the same way that un-collateralized loans do. In the worst
case, if customers default on their obligations, collateral will be
there as a back-up. More re-assuringly, the worst case is unlikely to
materialize since the fear of losing collateral should give customers
adequate incentives to repay loans. Investors will still worry about
currency risk, political risk, or macro-economic/environmental risk, but
investors can feel confident that the fundamental investment should be
unaffected by the institution's management structure, ownership rights,
or degree of leverage.

2.  \emph{The moderately-leveraged.} A second group of institutions
serves borrowers with less collateral than the first group. The
institution has to monitor and take other special efforts in order to
obtain ensure loan repayment. Given moral hazard, the more monitoring
that is required, the more the institution has to have incentives to
monitor. Here, the institution's owners have to be exposed to some of
the risk associated with the failure to monitor sufficiently. If they
didn't, potential investors would hesitate to invest. The microfinance
institution can thus achieve some leverage, but not as much as would be
possible when all loans are fully collateralized.

3.  \emph{The un-leveraged.} A third group of institutions serves
borrowers with even less collateral. Incentives are such that no
external leverage is possible at all. The institution is sustainable at
its existing scale. The institution mobilizes resources from its own
customers and proceeds on that basis---but commercial investors fail to
invest.

4.  \emph{The un-banked.} The fourth group of institutions are really
not institutions at all. They would fail to exist in a free market
system. Monitoring costs are so high that institutions would lose money
servicing customers. The customers may have worthy projects, but they
lack basic assets and are perceived to be so risky that they remain
un-banked. The potential customers would only be served if subsidies
emerged to launch the institutions.

We argue that there are two kinds of leverage that matter, and that this
recognition should push social investors either to focus on the
lower-end of markets that serve the poor or, for quite different
reasons, to focus on the higher-end. The first kind of leverage involves
bringing new pro-poor institutions into being but does not draw in
commercial capital. These are the latent institutions identified in the
fourth category above. Once in existence, the institutions can tap the
capital of their customers, and perhaps of others in the region, but
they are unlikely to be able to tap outside capital. By helping to bring
these institutions into existence, social investors can help put the
local capital to use and, in that way, multiply the power of their gift.

The second kind of leverage involves social investments that can ``crowd
in'' commercial capital to a maximal extent and the focus here is the
second category above. The first category of institutions, those that
are already highly leveraged in free markets, need little boost to reach
scale.

There is a middle ground in which little leverage is likely, captured by
the third category above, where institutions already exist in the free
market but they are not well-positioned to access commercial capital.
Social investments in the middle ground will have far less of a
multiplier effect than investments directed either up-market or
down-market. These institutions serve their customers well, and they
manage to tap their customers' deposits to maximize scale. But the
institutions are poorly-positioned to expand rapidly because of the
inherent riskiness of their portfolios and costs of their business, and,
if rapid expansion is the goal, social investors should look elsewhere.

Two points should be emphasized. First, social investment allocations
should also depend on the weights that investors attach to aiding the
various populations under consideration. Those weights may temper the
results here but will not affect the basic tradeoffs.

Second, we have assumed, without stressing it, that social investors are
not insiders in the institutions receiving funds. Investors have no
special knowledge and no special clout. In the language of corporate
finance, they constitute ``uninformed'' capital. The investors may be
Silicon Valley entrepreneurs, European religious groups, Japanese aid
agencies---each captured by the desire to maximize the social impact of
their resources, but retaining some distance from the ultimate
beneficiaries.

But imagine instead that investors have good information about the
diligence and management of the institution receiving the investment,
and that they have clout to affect the institutions' direction. If
investors are uninformed, they will have little means to instill good
incentives or mitigate moral hazard. If they are informed, though,
matters can be improved. In a world of limited information, uninformed
investors dominate the market (and they are vital in reaching scale),
but informed investors can have particular influence. The power of
informed social investors comes when they are able to combine capital
with knowledge, and in that way are able to attract ``uninformed''
investors. Incorporating this element into the framework will not change
the basic conclusions above, but it could open up new practical
strategies.

Little analytical has been written on social investment, or for that
matter on the economic incentives and constraints that guide
philanthropy. We have looked instead, partly out of necessity, to a very
different literature that shares important challenges and questions: the
economics of corporate finance. In building analytical links and
developing a simple framework for discussing new ideas, we hope to spur
others to take the analytical challenges seriously.

\hypertarget{appendix-1}{%
\section{\texorpdfstring{\\
Appendix 1}{ Appendix 1}}\label{appendix-1}}

Closed form solutions for \emph{N} and \emph{n} are obtained from the
system of equations

\((m,N) = A\)

\(N = \frac{(K - F)}{I^{m}}\)

Substitute
\(I_{m} = \frac{1}{\gamma}\left\lbrack \frac{\text{qm}}{\Delta} - \frac{F}{N} \right\rbrack\)from
into and the result into to solve for \emph{m}:

\[{N = \frac{(K - F)}{\frac{1}{\gamma}\left\lbrack \frac{\text{qm}}{\Delta} - \frac{F}{N} \right\rbrack}
}{N = \gamma\Delta N\frac{(K - F)}{\left\lbrack \text{Nqm} - \Delta F \right\rbrack}
}{m = \frac{1}{N}\frac{\Delta}{q}\left\lbrack \gamma K + (1 - \gamma)F \right\rbrack}\]

Substituting this back into to and re-arranging yields:

\[{(0) + \frac{F}{N} + m\left\lbrack 1 - \frac{\text{αp}}{\Delta} \right\rbrack = A
}{(0) + \frac{F}{N} + \left\lbrack \frac{1}{N}\frac{\Delta}{q}\left\lbrack \gamma K + (1 - \gamma)F \right\rbrack \right\rbrack\left\lbrack 1 - \frac{\text{αp}}{\Delta} \right\rbrack = A
}{F + \left\lbrack \frac{\Delta}{q}\left\lbrack \gamma K + (1 - \gamma)F \right\rbrack \right\rbrack\left\lbrack 1 - \frac{\text{αp}}{\Delta} \right\rbrack = N\left\lbrack A - (0) \right\rbrack}\]

\[\boxed{N = \frac{\Gamma}{\left\lbrack (0) - A \right\rbrack}}\]

Where
\(\Gamma = \frac{((1 - \alpha)p - q)}{q}\lbrack\gamma K - (1 - \gamma)F\rbrack - F\).

The closed form expressions for \emph{m} is then given by:

\[\boxed{\mspace{6mu} m = \left\lbrack (0) - A \right\rbrack \cdot \frac{\Delta}{q}\frac{\left\lbrack \beta K + (1 - \beta)F \right\rbrack}{\Gamma}}\]

Hence in equilibrium \emph{m} rises proportional to the gap between
\emph{{A}}(0) and \emph{A} -- poorer borrowers need more monitoring in
equilibrium. The number of borrowers reached by an MFI varies inversely
with the gap between \emph{{A}}(0) and \emph{A}

APPENDIX 2

(`Selfish equity' extension)

\textbf{2.A Selfish equity, no heterogeneity in each neighborhood}

In the current setup the rate of return on equity β was assumed to
equilibrate to γ, the same rate of return outside uninformed investors
expect on their investments. This assumption could be justified by
assuming that MFIs are run by social investors willing to accept such
return. This appendix examines the alternative scenario where it is
assume that the overall amount of `informed equity' (i.e. capital held
by individuals capable of managing a financial institution) in the
economy is limited, but that this informed equity can move across
`neighborhoods' in search of the highest return.

In order to isolate effects we make the additional simplifying
assumption that each neighborhood has the same finite number
\({\overline{N}}_{j}\)of would-be entrepreneurs with access to the same
lump sum investment project which gives gross return \emph{X} when it
succeeds and zero otherwise. In a world where investors are purely
profit driven, the first equity investors would at first want to direct
their scarce resources and talent to serving the wealthiest
neighborhoods for the simple reason that there is more project surplus,
and therefore economic rents to be made, where resources are not being
used up in costly monitoring. Competition for wealthiest clients should
drive down the rate of return in the wealthiest neighborhoods and lead
to a spillover of investments into less wealthy neighborhoods such that
in equilibrium the market rate of return on equity should be driven to
an economy-wide equilibrium level β which is the same in all
neighborhoods where clients are served.

To see how this equilibrium is reached, we first we allow the expected
rate of return in a given neighborhood to be set at arbitrary level β
and later solve for an equilibrium where β that clears the market for
\emph{K}.

When there are fixed costs to setting up an institution one should
really have to be worried about strategic behavior and timing. In the
math below we have included fixed costs per borrower in most expressions
to allow for the possibility of analyzing this at some later date, but
many of the arguments really implicitly assume \emph{F}=0.

The main result of interest is that the pursuit of higher profits
changes the conditions of access and will in general lead to
`over-monitoring' in the market. To see this note the rate of return on
equity is:

Or put only slightly differently, the expected earnings to the equity
investor/monitor from an incentive-compatible contract must exactly
equal the opportunity cost of those funds \emph{I\textsubscript{m}} in
other equity investments:

\(\beta I_{m} = \left\lbrack \frac{q}{\Delta}m - \frac{F}{N} \right\rbrack\)

Since the rate of return β demanded by for-profit equity investors will
in general exceed the rate of return γ that equity owners obtained in
the earlier analyzed `most-generous social investor' case, borrowers
(who were contributing as much collateral as they had in that earlier
scenario) must now be more heavily monitored if incentives are to be
kept in place for them to credibly commit to these higher levels of
expected repayment. To see this formally note that the contract must now
divide expected project claims \emph{pX} between sufficient expected
payments for the uninformed investor (who puts up\(I - I_{m}\)) to
participate, sufficient expected returns for the borrower to be
diligent, and expected payments \(pR_{m} = \frac{p}{\Delta}m\) to the
monitor/equity owner:

\(pX = \gamma(I - I_{m}) + \left\lbrack \frac{pB(m)}{\Delta} - A \right\rbrack + \left\lbrack \frac{p}{\Delta}m \right\rbrack\)

Re-arranging to solve for A as before we get:

where\((m,\gamma)\) was the minimum collateral requirement in the
`generous social investors' case where the rate of return on equity was
assumed to be β=γ (i.e. equity investors accepted the same rate of
return as uninformed investors).

Substituting for \emph{I\textsubscript{m}} from above allows us to
further re-arrange:

Which is all a lengthy way to say
\((m,\gamma) = (m,\gamma) + (\beta - \gamma)I_{m}\)

**** COMMENT ***

The following figure illustrates how the monitoring intensity must rise.
The lower line is the minimum collateral requirement for the social
investor case (i.e. β=γ=1.1), while the top line is the minimum
collateral requirement is 20\% higher, or γ=1.32). A borrower with say
\(A_{j} = 100\)would obtain a loan with (approximately) 7.5 units of
required monitoring in the first place, while the same borrower would
require 12.9 units of monitoring

\includegraphics[width=3.81806in,height=2.375in]{media/image10.wmf}

*** END COMMENT ***

\textbf{`Closing' the model }

Suppose that there are \emph{j} = \emph{1...J} neighborhoods with
\({\overline{N}}_{j} = N\) potential entrepreneur/borrowers in each
neighborhood, who differ only in terms of their collateralizable wealth
\emph{A}\textsubscript{j}.

Equity investor funds will flow to the wealthiest neighborhoods until
all borrowers are served. Competition will however drive down the rate
of return in this neighborhood making it attractive for funds to flow to
the next richest neighborhood, and so on.

The logic runs loosely like this. Suppose there were just two
neighborhoods, one rich, one poor. Suppose K is in very scarce supply.
It's obvious that these equity investors will setup a bank in the richer
neighborhood since that is where they can leverage more outside
investment, and hence earn a higher return on equity. These
profit-seeking investor would like to be as highly leveraged as
possible, but competition amongst them to attract rich borrowers will
force them to be less highly leveraged

Any remaining funds will flow to the next richest neighborhood until all
borrowers are served, and so on, flowing to the next richest
neighborhood each time. In the last neighborhood served

Conceptually we can solve for demand for K in each neighborhood as
follows. Take neighborhood \emph{A\textsubscript{j}} . From

How does shape change if β\textgreater{}γ: If the MFI demands a higher
return on equity β (perhaps because monitoring equity capital is scarce
in the economy -- we can endogenize this rate of return by specifying a
fixed quantity of monitoring equity in the economy) the higher share of
project surplus devoted to equity payments shifts the thresholds and
reduces access.

\includegraphics[width=6.19792in,height=2.52083in]{media/image11.wmf}

Access is reduced (not surprisingly since the MFI is taking more surplus
off the table). Note that there is also now a higher level of monitoring
at each level of \emph{A}. This is because if the MFI is to take more
money off the table, the borrowers must have incentives to be able to
commit to deliver it. Or put differently, with less money for them to
take home when things succeed their incentive to be diligent would be
diluted, unless monitoring can be stepped up.

THERE ARE TWO EXPRESSIONS IN TWO UNKNOWS (m, N):

\[N = \frac{K - F}{I_{m}}\]

Substituting in
\(I_{m} = \frac{1}{\beta}\left\lbrack \frac{p \cdot m}{\Delta} - m - \frac{F}{N} \right\rbrack\)
these can be rewritten in terms of just N and m:

Note that if β=γ the expression reduces to:

\(m = \left\lbrack \frac{(p - q)}{(\alpha - 1)p + q} \right\rbrack\left\lbrack \left\lbrack {}_{0} - A \right\rbrack + \frac{F}{N} \right\rbrack\)

AND

\(\boxed{m = \left( \frac{\beta K + (1 - \beta)F}{N} \right)\left( \frac{p - q}{q} \right)}\)

Solving for N

\[{\beta K + (1 - \beta)F = N\left\lbrack \frac{\text{βq}}{(\alpha - 1)\beta p + \gamma q} \right\rbrack\left\lbrack \left\lbrack {}_{0} - A \right\rbrack + \frac{\gamma}{\beta}\frac{F}{N} \right\rbrack
}{\left\lbrack \beta K + (1 - \beta)F \right\rbrack\left\lbrack \frac{(\alpha - 1)\beta p + \gamma q}{\text{βq}} \right\rbrack = N\left\lbrack {}_{0} - A \right\rbrack + \frac{\gamma}{\beta}F
}{N = \frac{1}{\left\lbrack {}_{0} - A \right\rbrack}\left\lbrack \left\lbrack \beta K + (1 - \beta)F \right\rbrack\left\lbrack \frac{(\alpha - 1)\beta p + \gamma q}{\text{βq}} \right\rbrack - \frac{\gamma}{\beta}F \right\rbrack
}\boxed{N = \frac{\Gamma(\beta,\gamma)}{\left\lbrack {}_{0} - A \right\rbrack}}\]

From here I should be easy to show that, all else equal, an increase β
lowers equilibrium N and raise equilibrium m.

Borrower's participation constraint

\(\text{pX} - \text{γI} - \left\lbrack \frac{\text{pm}}{\Delta} - \gamma I_{m} \right\rbrack \geq 0\)

Which can b rewritten

\[{\text{pX} - \text{γI} - \left\lbrack \frac{\text{pm}}{\Delta} - \gamma I_{m} \right\rbrack \geq 0
}{\text{pX} - \text{γI} - \frac{\text{pm}}{\Delta} + \frac{\gamma}{\beta}\left\lbrack \frac{\text{pm}}{\Delta} - m - \frac{F}{N} \right\rbrack \geq 0
}{\text{pX} - \text{γI} - m\left( \frac{\text{βp} - \text{γq}}{\beta\Delta} \right) - \frac{\gamma}{\beta}\frac{F}{N} \geq 0}\]

When if β=γ the expression reduces to:
\(= pX - \text{γI} - \frac{F}{N} - m\)

\includegraphics[width=6.03125in,height=2.65625in]{media/image13.wmf}

\hypertarget{section}{%
\section{}\label{section}}

\hypertarget{section-1}{%
\section{\texorpdfstring{\\
}{ }}\label{section-1}}

\hypertarget{bibliography}{%
\section{Bibliography}\label{bibliography}}

Abrams, J. and D. von Stauffenberg (2007). Role Reversal: Are Public
Development Institutions Crowding Out Private Investment in
Microfinance? Microrate, MFI Insights, February 2007.
http://www.microrate.com/pdf/rolereversal.pdf

Armendariz de Aghion, B. and J. Morduch (2005). \emph{The Economics of
Microfinance}, MIT Press.

Barrés, Isabelle (2005) ``Bulletin Highlights-Supply of Funding,''
\emph{Microbanking Bulletin}, August 2005, p.53.

Becchetti, L. et al. (2005) , ``A matching of two promises: Microfinance
and Social Responsibility'', Dipartimento di Economia e Istituzioni,
Università degli Studi di Roma "Tor Vergata".

Berlin, M. (1998). "That Thing Venture Capitalists Do." \emph{Business
Review}, Federal Reserve Bank of Philadelphia: 15-26.

Bardhan, Pranab, Samuel Bowles and Herbert Gintis (2000), ``Wealth
Inequality, Wealth Constraints, and Economic Performance,'' in Atkinson
and Bourguignon eds. \emph{Handbook of Economic Growth,} North-Holland.

Chowdhry, B. et al. (2005), ``Pricing Microfinance Loans and Loan
Guarantees using Biased Loan write-off Data'', Manuscript.

CGAP (2004), ``Building Inclusive Financial Systems: Donor Guidelines on
Good Practice in Microfinance''.
http://cgap.org/docs/donorguidelines.pdf

Conning, J. (1999) ``Outreach, Sustainability and leverage in monitored
and peer-monitored lending,'' \emph{Journal of Development Economics},
Vol. 60, 51-77.

Conning, J. (2006). ``Monitoring by Delegates or by Peers? Joint
Liability Loans under Moral Hazard,'' Hunter College, Department of
Economics, working paper.
\href{http://ideas.repec.org/p/htr/hcecon/407.html}{ideas.repec.org/p/htr/hcecon/407.html}

Conning, J. and C. Udry (2007). ``Rural Financial Markets in Developing
Countries.'' Chapter 56 of \emph{Handbook of Agricultural Economics},
Vol. 3. R.E. Evenson, P. Pingali, and T. P. Schultz (ed.). Amsterdam:
Elsevier.

Cull, R. A. Demirgüç-Kunt, and J. Morduch (2008) ``Microfinance Meets
the Market,'' \emph{Journal of Economic Perspectives}, forthcoming.

Dewatripont, M. and J. Tirole (1994). \emph{The Prudential Regulation of
Banks}. Cambridge, Mass., MIT Press.

Dewatripont, M. and G. Roland (2000). "Soft Budget Constraints,
Transition, and Financial Systems." \emph{Journal of Institutional and
Theoretical Economics} 156(1): 245-60.

Diamond, D. W. (1984). "Financial Intermediation and Delegated
Monitoring." Review of Economic Studies 51(3): 393-414.

Hansmann, Henry (1996), \emph{The Ownership of Enterprise. Cambridge},
Harvard University Press.

Heinen, Erik (2005) ``The MFI as a Borrower: Institutional
Characteristics and MFI Performance,'' Microbanking Bulletin, August,
p.33.

Innes, R. D. (1990). "Limited Liability and Incentive Contracting with
Ex-ante Action Choices." \emph{Journal of Economic Theory} 52(1): 45-67.

Lauer, Kate. (2008). ``Transforming NGO MFIs: Critical Ownership Issues
to Consider.'' \emph{CGAP Occasional Paper No. 13}.

Ledgerwood, J. (1998). Microfinance Handbook : An institutional and
financial perspective, World Bank Publications.

Li, Wenli (1998) ``Government Loan, Guarantee, and Grant Programs: An
Evaluation,'' \emph{Federal Reserve Bank of Richmond Economic Quarterly}
84, no. 4.

Marble, Melinda (1989), \emph{Social Investment and Private
Foundations}, Council on Foundations, Washington, DC.

Mody, A. and Patro, D. (1996), ``Methods of Loan Guarantee Valuation and
Accounting,'' Manuscript.

Shiller, R. J. 2008. The Subprime Solution: How Today's Global Financial
Crisis Happened and what to Do about it. Princeton University Press.

Tirole, Jean (2005), \emph{The Theory of Corporate Finance}. Princeton:
Princeton University Press.

Tirole, J. (2001). ``Corporate Governance,'' \emph{Econometrica} 69 (1):
1-35.

Townsend, R. M. (1979). "Optimal Contracts and Competitive Markets with
Costly State Verification." \emph{Journal of Economic Theory} 21(2):
265-93.

Weinsberg, Jacob (2006). ``The Philanthropist's Handbook: How
billionaires give their money away.'' \emph{Slate} (on-line). Posted
Wednesday, Nov. 15, 2006, at 3:30 PM ET:
http://www.slate.com/id/2153857/.

  

SOME FIGURES FROM THE SLIDE PRESENTATION:

\textbf{`Profitability' versus non-commercial funding.} Profitability is
measured by the financial sustainability ratio. The financial
self-sufficiency ratio is adjusted financial revenue divided by the sum
of adjusted financial expenses, adjusted net loan loss provision
expenses, and adjusted operating expenses. It indicates the
institution's ability to operate without ongoing subsidy, including soft
loans and grants. The definition is from \emph{MicroBanking Bulletin}
(2005), p. 57.

\includegraphics[width=5.21389in,height=3.81806in]{media/image14.emf}

\includegraphics[width=4.98958in,height=3.625in]{media/image15.wmf}

\textbf{The composition of costs as a share of the average gross loan
portfolio.} Data on costs are not disaggregated by activity, so the
analysis includes only those institutions whose revenue from lending
(interest, fees, and commissions) is greater than or equal to 80 percent
of total revenues. In restricting ourselves to the subset of
microfinance institutions that are most focused on lending, we have
greater confidence in ascribing all of their costs to lending.

\includegraphics[width=5.16667in,height=3.76042in]{media/image16.wmf}

\textbf{Interest rates rise with costs.} The ``premium'' is the excess
of the microlender's average interest rate charged to borrowers over the
International Monetary Fund's inter-bank ``lending interest rate'' that
banks in the given countries charge to prime customers (from IMF
International Financial Statistics).

TABLES FROM THE PAPER

Table : Numerical Example

\begin{longtable}[]{@{}lll@{}}
\toprule
\textbf{Variable or Parameter} & \textbf{Value} &
\textbf{Description}\tabularnewline
\midrule
\endhead
\emph{I} & 100 & Lump sum investment\tabularnewline
\emph{X} & 200 & Successful Project outcome (zero for
failure)\tabularnewline
\emph{p} & 0.90 & Probability of Project Success when
diligent\tabularnewline
\emph{q} & 0.75 & Probability of Project Success when not
diligent\tabularnewline
\(\gamma\) & 1.20 & Lender's cost of funds (one plus interest
rate)\tabularnewline
\emph{F} & 1000 & Fixed cost to setup or operate an MFI\tabularnewline
\emph{K-F} & 11000 & Initial capital of the MFI\tabularnewline
\(B_{0}\) & 30 & Opportunity cost of diligence at zero
monitoring\tabularnewline
\(\alpha\) & 0.5 & Monitoring effectiveness parameter\tabularnewline
\(B(m)\) & \(B_{0} - \alpha \cdot m\) & Monitoring
technology\tabularnewline
\(N_{0}\) & \(\frac{K - F}{I} = 100\) & Max number of borrowers reached
(without leverage)\tabularnewline
\(AC_{0}\) & \(\frac{F}{N_{0}} = 10\) & Average fixed cost per borrower
(without leverage)\tabularnewline
\bottomrule
\end{longtable}

Figure : How MFIs fund themselves

\includegraphics[width=6.26597in,height=3.39097in]{media/image17.emf}

\includegraphics[width=6.27083in,height=2.57847in]{media/image18.emf}

\includegraphics[width=5.58333in,height=1.34375in]{media/image19.png}

Figure : Collateral and monitoring

\includegraphics[width=8.55387in,height=9.35517in]{media/image20.emf}

\includegraphics[width=6.26042in,height=7.67222in]{media/image24.png}

\includegraphics[width=6.27083in,height=5.05208in]{media/image25.png}

\includegraphics[width=6.26597in,height=3.625in]{media/image26.png}

\end{document}
